\documentclass[11pt,compress,t,notes=noshow, xcolor=table]{beamer}

\input{../../style/preamble}
\input{../../latex-math/basic-math}
\input{../../latex-math/basic-ml}
\input{../../latex-math/optim-basics}

% Lipschitz constant of the Hessian (Boyd's L; plain L is already the smoothness
% constant in this lecture, see 01-foundations/03-convex.tex)
\newcommand{\LH}{L_{\H}}

\title{Optimization in Machine Learning}

\begin{document}

\titlemeta{
Second order methods
}{
Deep Dive: Convergence of Newton-Raphson
}{
figure/NR_phases.png
}{
\item Assumptions
\item Damped phase
\item Quadratically convergent phase
\item Complexity bound
\vspace{0.5cm}
\item[] \hspace{-0.7cm}Slides based on\\
\hspace{-0.7cm}\furtherreading{BOYD2004} (Ch. 9.5.3)
}

\begin{framei}{Setting}
\item \textbf{Recall:}
\begin{itemizeM}
\item with full steps ($\alpha^{[t]} = 1$), NR may overshoot/diverge far from
$\xvs$ $\Rightarrow$ damping is necessary
\item many objective functions are only locally quadratic
\end{itemizeM}
\item In this deep dive, we show that damped NR has two phases:
\begin{itemizeM}
\item \textbf{Damped phase:} quadratic model unreliable, $\alpha^{[t]} < 1$ possible; $f$ decreases by a
\textit{constant} $\delta$ per step
\item \textbf{Pure Newton phase:} quadratic model accurate, $\alpha^{[t]} = 1$ always accepted; \textit{quadratic} convergence
\end{itemizeM}
\splitV[0.62]{
\imageC[0.85]{figure/NR_phases.png}
}{
{\footnotesize \textbf{Example 2} (cont'd):
$$f(\xv) = (1 + x_1^2)^{1/2} + \frac{x_2^2}{2}$$
from $\xv^{[0]} = (3, 2)^\top$,\\
with $\gamma_1 = 0.1$, $\tau = 0.7$}
}
\end{framei}


\begin{framei}[fs=small]{Setting}
\item \textbf{Assumptions} on the sublevel set $\S = \{\xv : f(\xv) \le f(\xv^{[0]})\}$:
\begin{itemizeM}
\item $f \in \CC{2}$, $m$-strongly convex ($\H(\xv) \succeq m \id$) and $L$-smooth ($\H(\xv) \preceq L \id$)
\item $\H$ is \textbf{Lipschitz} with constant $\LH$ ($\rightarrow$ bound on the \textit{3rd} derivative):
$$\|\H(\xv) - \H(\xtil)\|_2 \le \LH \|\xv - \xtil\| \qquad \text{for all } \xv, \xtil \in \S$$
Intuition: $\LH$ measures how well the quadratic model approximates $f$ ($\LH = 0$ for quadratic $f$)
\end{itemizeM}
\item Damped NR with backtracking: $\alpha_{\text{init}} = 1$, $\gamma_1 \in (0, 1/2)$, $\tau \in (0, 1)$
\begin{framed}
\textbf{Theorem \textnormal{(Convergence of damped NR)}.}
With the scaled gradient norm $\theta^{[t]} = \frac{\LH}{2m^2}\|\nabla f(\xv^{[t]})\|$, there exist
$\eta \in (0, m^2/\LH]$ and $\delta > 0$ such that
\begin{align*}
\|\nabla f(\xv^{[t]})\| \ge \eta \; &\Rightarrow \; f(\xv^{[t+1]}) - f(\xv^{[t]}) \le - \delta
&& \textit{(damped)} \\
\|\nabla f(\xv^{[t]})\| < \eta \; &\Rightarrow \; \alpha^{[t]} = 1 \, \text{ and } \,
\theta^{[t+1]} \le \big( \theta^{[t]} \big)^2 && \textit{(pure)}
\end{align*}
\end{framed}
\end{framei}

% \begin{framei}{Outline}
% \item We start by analyzing the implications of the theorem (main convergence result)
% \item We then actually prove the theorem
% \begin{itemizeM}
% \item \textbf{Damped phase:} We show that each step reduces $f$ by at least a constant $\delta$ and show how $\delta$ looks like
% \item \textbf{Pure phase:} We show that $\alpha^{[t]} = 1$ is always accepted and that the scaled gradient norm $\theta^{[t]}$ decreases quadratically
% \end{itemizeM}
% \end{framei}

\begin{framei}[fs=small]{What the theorem implies: pure phase}
\item [] Before we actually prove the theorem, we analyze its implications for the convergence of damped NR.
\item \textbf{Message 1:} once NR enters the pure phase, it stays there
\begin{itemizeM}
  \item suppose we are in the pure phase at step $t$, i.e. $\|\nabla f(\xv^{[t]})\| < \eta$
  \item since $\eta \le m^2/\LH$ by def., the scaled gradient norm is below $\tfrac{1}{2}$:
  $$\theta^{[t]} = \frac{\LH}{2m^2}\|\nabla f(\xv^{[t]})\|
  < \frac{\LH}{2m^2} \cdot \frac{m^2}{\LH} = \frac{1}{2}$$
  \item with this and the theorem:
  $$\theta^{[t+1]} \le \big( \theta^{[t]} \big)^2 = \theta^{[t]} \cdot \theta^{[t]}
  < \tfrac{1}{2} \theta^{[t]}$$
  $$\Rightarrow \|\nabla f(\xv^{[t+1]})\| < \tfrac{1}{2} \|\nabla f(\xv^{[t]})\| < \eta$$
  \item i.e., the condition is also satisfied at step $t+1$ and by recursion at all future steps $l \geq t$, implying (by the theorem):
  $$\alpha^{[l]} = 1 \quad \text{ and } \quad \theta^{[l+1]} \le \big( \theta^{[l]} \big)^2$$
\end{itemizeM}
\end{framei}

\begin{framei}[fs=small]{What the theorem implies: pure phase}
\item \textbf{Message 2:} convergence is extremely rapid in the pure phase
\begin{itemizeM}
\item $\theta^{[t+1]} \le \big( \theta^{[t]} \big)^2$ is exactly \textbf{quadratic convergence}
\item Applying it recursively, for $l \geq t$:
$$\theta^{[l]} \le \big( \theta^{[t]} \big)^{2^{l-t}} \leq \left(\frac{1}{2}\right)^{2^{l-t}}$$
\item Strong convexity gives $f(\xv) - f(\xvs) \le \frac{1}{2m}\|\nabla f(\xv)\|^2$ and thus:
$$f(\xv^{[l]}) - f(\xvs) \le \eps_0 \big( \theta^{[l]} \big)^2
\le \eps_0 \left(\frac{1}{2}\right)^{2^{l-t+1}}, \qquad \eps_0 = \frac{2m^3}{\LH^2}$$
\item Inverting this for the step count yields:\footnote{\footnotesize Exactly
$\log_2 \log_2 (\eps_0/\eps) - 1$ steps; the $-1$ is dropped for readability.}
\begin{framed}
$$f(\xv^{[l]}) - f(\xvs) \le \eps \quad \text{after at most} \quad \log_2 \log_2 (\eps_0/\eps)
\; \text{ pure steps}$$
\end{framed}
\item[$\Rightarrow$] \textbf{Note}: at most 6 pure steps give an accuracy of
$\eps \approx 5 \cdot 10^{-20} \, \eps_0$
\end{itemizeM}
\end{framei}

\begin{framei}[fs=small]{What the theorem implies: damped phase \& final result}
\item \textbf{Message 3:} the number of steps in the damped phase is bounded
\begin{itemizeM}
\item Every damped step reduces $f$ by at least $\delta$, so after $T$ damped steps
$f(\xv^{[T]}) \le f(\xv^{[0]}) - T \delta$
\item Since $f$ cannot decrease below $f(\xvs)$:
\begin{framed}
$$\#\{\text{damped steps}\} \le \frac{f(\xv^{[0]}) - f(\xvs)}{\delta}$$
\end{framed}
\item[$\Rightarrow$] finite, but expensive phase; grows with the initial gap
\end{itemizeM}
\item \textbf{Final result} (both phases together): to reach $f(\xv^{[t]}) - f(\xvs) \le \eps$:
\begin{framed}
$$\#\text{iterations} \;\; \le \;\;
\underbrace{\frac{f(\xv^{[0]}) - f(\xvs)}{\delta}}_{\text{damped}} \; + \;
\underbrace{\log_2 \log_2 (\eps_0/\eps)}_{\text{pure}, \; \approx \, 6}$$
\end{framed}
\end{framei}


\begin{framev}[fs=small]{Proof strategy}
\textbf{Note:} we have now seen the implications of the theorem, but have not yet proven it, i.p. the existence of $\eta$ and $\delta$. Our proof strategy is as follows:
\begin{enumerate}
\item \textbf{Tool:} the \textbf{Newton decrement} $\lambda(\xv)$ as progress measure
\item \textbf{Damped phase:} $L$-smoothness bounds the step size from below
\begin{framed}
$$\alpha^{[t]} \ge \tau \frac{m}{L}
\quad \Longrightarrow \quad
\delta = \gamma_1 \tau \eta^2 \frac{m}{L^2}$$
\end{framed}
\item \textbf{Pure phase:} Lipschitz $\H$ keeps the quadratic model accurate enough that
\begin{framed}
$$\alpha^{[t]} = 1 \text{ is accepted}
\quad \text{and} \quad
\theta^{[t+1]} \le \big( \theta^{[t]} \big)^2$$
\end{framed}
\end{enumerate}
\end{framev}


\begin{framev}[align=top]{Discussion}
  tbd: connect to statements from first deck (locally, quadratic convergence, etc.) and discuss what would happen for non-damped NR, maybe also the role of the assumptions
\end{framev}

\endlecture
\end{document}
